\documentclass[11pt]{article}
\usepackage[margin=1in]{geometry}
\usepackage{graphicx}
\usepackage{hyperref}
\usepackage{algorithm,algorithmic,caption}
\usepackage{amsmath}
\usepackage[T1]{fontenc}
\usepackage{amsfonts}
\usepackage{enumitem}
\usepackage{natbib}
\usepackage{listings}
\usepackage{xcolor}
\usepackage{tcolorbox}
\usepackage{titlesec}
\usepackage{fancyhdr}

% Modern color palette
\definecolor{primaryblue}{RGB}{37,99,235}
\definecolor{lightblue}{RGB}{59,130,246}
\definecolor{darkblue}{RGB}{30,64,175}
\definecolor{accentgray}{RGB}{100,116,139}
\definecolor{lightgray}{RGB}{248,250,252}

% Hyperlink setup
\hypersetup{
    colorlinks=true,
    linkcolor=primaryblue,
    urlcolor=lightblue,
    citecolor=darkblue
}

% Section styling
\titleformat{\section}
  {\color{primaryblue}\Large\bfseries}
  {\thesection}{1em}{}
  [\titlerule]

\titleformat{\subsection}
  {\color{lightblue}\large\bfseries}
  {\thesubsection}{1em}{}

% Page styling
\pagestyle{fancy}
\fancyhf{}
\fancyhead[L]{\color{accentgray}\small Graphs in Machine Learning}
\fancyhead[R]{\color{accentgray}\small TD3: Graph Neural Networks}
\fancyfoot[C]{\color{accentgray}\thepage}
\renewcommand{\headrulewidth}{0.5pt}
\renewcommand{\headrule}{\hbox to\headwidth{\color{lightblue}\leaders\hrule height \headrulewidth\hfill}}

% Paragraph settings
\setlength\parindent{0pt}
\setlength{\parskip}{0.5\baselineskip}

% Custom boxes
\newtcolorbox{infobox}{
  colback=lightgray,
  colframe=lightblue,
  fonttitle=\bfseries,
  title=Key Information,
  arc=3mm,
  boxrule=1.5pt
}

\newtcolorbox{questionbox}{
  colback=lightgray,
  colframe=primaryblue,
  fonttitle=\bfseries,
  title=Questions,
  arc=3mm,
  boxrule=1.5pt
}

\begin{document}

% Title page
\begin{center}
    \includegraphics[width=0.4\textwidth]{../../../../../images/mva/peer_communities2.png}\\[1cm]

    {\Huge\color{primaryblue}\textbf{Practical Session 3}}\\[0.5cm]
    {\LARGE\color{darkblue}Graph Neural Networks}\\[1cm]

    {\large\color{accentgray}
    Graphs in Machine Learning\\
    MVA Master Program -- ENS Paris-Saclay}\\[1.5cm]

    {\color{accentgray}
    \textit{michal.valko@inria.fr}\\
    \textit{daniele.calandriello@inria.fr}\\
    \textit{omar.darwiche-domingues@inria.fr}\\
    \textit{pierre.perrault@inria.fr}}
\end{center}

\vfill

\begin{tcolorbox}[colback=lightgray,colframe=lightblue,arc=2mm,boxrule=1pt]
\textbf{About this Session}\\[0.3cm]
Complete the exercises and submit a written report.
\end{tcolorbox}

\newpage

\section{Neural Relational Inference}

This practical session is based on the paper
\href{https://arxiv.org/pdf/1802.04687.pdf}{\textit{Neural Relational Inference for Interacting Systems}}
by Kipf et al., 2018.

We will use material provided by Marc Lelarge and Timothée Lacroix:

\begin{center}
\url{https://github.com/timlacroix/nri_practical_session}
\end{center}

\subsection{Motivation and Problem Formulation}

A wide range of dynamical systems can be seen as a group of interacting components.
For example, we can think of a set of 2-dimensional particles coupled by springs.
Assume we are given only a set of trajectories of such an interacting dynamical system.
How can we learn its dynamical model in an unsupervised way?

\begin{infobox}
\textbf{Formal Setup:}

We are given as input trajectories of $N$ objects, each trajectory has length $T$.

\begin{itemize}[leftmargin=*]
    \item Each object $i$, for $i=1,\ldots, N$, is represented by vertex $v_i$
    \item Let $\mathbf{x}_i^t$ be the feature vector of object $i$ at time $t$ (e.g., position and velocity) with dimension $D$
    \item Let $\mathbf{x}^t = \{\mathbf{x}_1^t, \ldots, \mathbf{x}_N^t\}$ be features of all $N$ objects at time $t$
    \item Let $\mathbf{x}_i = (\mathbf{x}_i^1, \ldots, \mathbf{x}_i^T)$ be trajectory of object $i$
    \item Input data: 3D array $\mathbf{x}$ of shape $N\times T \times D$, denoted by $\mathbf{x} = (\mathbf{x}^1, \ldots, \mathbf{x}^T)$
\end{itemize}

The element $\mathbf{x}_{i,t,d}$ is the $d$-th component of the feature vector of object $i$ at time $t$.
\end{infobox}

We assume the dynamics can be modeled by a graph neural network (GNN) given an unknown graph $\mathbf{z}$ where $\mathbf{z}_{i,j}$ represents the discrete edge type between objects $v_i$ and $v_j$.

\subsection{Learning Objectives}

In this context, we want to learn simultaneously:

\begin{enumerate}
	\item \textbf{Edge type estimation}: The edge types $\mathbf{z}_{i,j}$
	\item \textbf{Future state prediction}: A model that, for any time $t$, takes $\mathbf{x}^t$ as input and predicts $\mathbf{x}^{t+1}$ as output
\end{enumerate}

\subsection{Model Architecture}

The Neural Relational Inference (NRI) model consists of:

\begin{infobox}
\textbf{NRI Components:}

\begin{itemize}[leftmargin=*]
	\item \textbf{Encoder:} Uses trajectories $\mathbf{x} = (\mathbf{x}^1, \ldots, \mathbf{x}^T)$ to infer pairwise interaction vectors $\mathbf{z}_{i, j} \in \mathbb{R}^K$ for $i, j$ in $\{1, \ldots, N \}$, where $K$ is the number of \emph{edge types}

	\item \textbf{Decoder:} Takes $\mathbf{x}^t$ and $\mathbf{z} = \{\mathbf{z}_{i, j}\}_{i,j}$ as input to infer $\mathbf{x}^{t+1}$
\end{itemize}
\end{infobox}

Both the encoder and the decoder are implemented using graph neural networks.
For more details, read Section 3 of the paper
\href{https://arxiv.org/pdf/1802.04687.pdf}{here}.

\section{Exercises}

Complete the code in the following notebook:

\begin{center}
\url{https://github.com/timlacroix/nri_practical_session/blob/master/NRI_student.ipynb}
\end{center}

and answer the questions below in your report.

\begin{tcolorbox}[colback=lightgray,colframe=darkblue,arc=2mm,boxrule=1pt]
\textbf{Note:} For the report, no code submission is required. The Github repository
contains a \texttt{solutions} folder, which you are allowed to use to complete the notebook.
\end{tcolorbox}

\begin{questionbox}
\begin{enumerate}[leftmargin=*]
	\item Explain what are the edge types $\mathbf{z}_{i,j}$.

	\item In the NRI model, explain how the encoder and the decoder work.

	\item Explain the LSTM baseline used for joint trajectory prediction.
	      Why is it important to have a ``burn-in'' phase?

	\item Consider the training of the LSTM baseline. Notice that the negative log-likelihood
	      is lower after the burn-in than before. Why is this surprising? Why is this happening?

	\item Consider the problem of trajectory prediction. What are the advantages of the NRI model
	      with respect to the LSTM baseline?

	\item Consider the training of the NRI model. What do you notice about the edge accuracy
	      during training? Why is this surprising?

	\item What do you expect to happen with the NRI model when there is no interaction
	      between the objects?
\end{enumerate}
\end{questionbox}

\end{document}
